\chapter{课堂练习及结果}
\label{cp:introduction}

\section{第 9 题 从源码构建并在干净环境安装 Wheel}

\subsection{实验内容}

如图\ref{fig:w3目录}所示，在课程仓库中创建 w3 实验目录并进入，第三周的四个实验均在 w3 目录下完成。

\begin{figure}[H]
    \centering
    \includegraphics[width=0.8\textwidth]{截屏2026-09-09 10.45.13.png}
    \caption{创建 w3 实验目录}
    \label{fig:w3目录}
\end{figure}

由于本机没有预装课程要求的构建工具，需要先准备课程环境。如图\ref{fig:课程环境}所示，先检查 Python 版本（3.14.7），再使用 \verb|python3 -m venv| 命令创建课程虚拟环境 \verb|course_env|；如图\ref{fig:安装课程工具}所示，在 \verb|course_env| 中执行 \verb|pip install build pytest| 安装构建和测试工具，为后续构建 wheel 和运行测试做准备。

\begin{figure}[H]
    \centering
    \includegraphics[width=0.8\textwidth]{截屏2026-09-09 10.53.22.png}
    \caption{检查 Python 版本并创建课程虚拟环境 course\_env}
    \label{fig:课程环境}
\end{figure}

\begin{figure}[H]
    \centering
    \includegraphics[width=0.8\textwidth]{截屏2026-09-09 10.54.09.png}
    \caption{在 course\_env 中安装 build 与 pytest}
    \label{fig:安装课程工具}
\end{figure}

如图\ref{fig:创建greetlab包}所示，在 q09 中创建 \verb|src/greetlab/__init__.py|、cli.py 和 pyproject.toml。其中 \verb|[project]| 段的 name 为 \texttt{greetlab-25020007076}（把题目模板中的"学号"替换为自己的学号）、version 为 0.1.0；\verb|[project.scripts]| 段把命令行入口 \verb|sdt-greet| 指向 \verb|greetlab.cli| 的 main 函数。

\begin{figure}[H]
    \centering
    \includegraphics[width=0.8\textwidth]{截屏2026-09-09 10.55.55.png}
    \caption{创建 cli.py 与 pyproject.toml}
    \label{fig:创建greetlab包}
\end{figure}

如图\ref{fig:构建wheel}所示，激活 \verb|course_env| 后运行 \verb|python -m build|，成功生成 \texttt{greetlab\_25020007076-0.1.0.tar.gz} 和 \texttt{greetlab\_25020007076-0.1.0-py3-none-any.whl} 两个分发产物。

\begin{figure}[H]
    \centering
    \includegraphics[width=0.8\textwidth]{截屏2026-09-09 10.56.57.png}
    \caption{运行 python -m build 生成 wheel}
    \label{fig:构建wheel}
\end{figure}

如图\ref{fig:干净环境安装}所示，先用 \verb|ls dist| 查看构建产物，再使用 \verb|python -m venv venv| 新建一个干净虚拟环境，最后用 \verb|venv/bin/pip install dist/greetlab_25020007076-0.1.0-py3-none-any.whl| 只从生成的 wheel 安装，输出 \texttt{Successfully installed greetlab-25020007076-0.1.0}，全程没有从镜像源下载安装该包。

\begin{figure}[H]
    \centering
    \includegraphics[width=0.8\textwidth]{截屏2026-09-09 10.57.59.png}
    \caption{新建干净虚拟环境并仅从 wheel 安装}
    \label{fig:干净环境安装}
\end{figure}

如图\ref{fig:sdt-greet运行}所示，切换到 q09 之外（w3 目录）运行 \verb|sdt-greet --name 25020007076|，输出 \texttt{Hello, 25020007076!}，说明 wheel 安装正确、命令行入口可用。

\begin{figure}[H]
    \centering
    \includegraphics[width=0.8\textwidth]{截屏2026-09-09 10.58.49.png}
    \caption{在 q09 之外运行 sdt-greet}
    \label{fig:sdt-greet运行}
\end{figure}

\subsection{实验结果}

第九题完成了从源码构建 wheel、在干净虚拟环境中仅从 wheel 安装、并在 q09 之外运行命令行工具验证输出的完整流程，验证了 Python 打包（Packaging）与安装（Shipping）的整条链路。

% TODO: 完成 w3 目录的 git 提交与推送后，将 commit/push 记录截图放入 images/ 并补充如下内容
% 第九题的 github 的 commit 记录和 push 记录如图\ref{fig:github记录Q9}所示。
% \begin{figure}[H]
%     \centering
%     \includegraphics[width=0.8\textwidth]{github记录Q9.png}
%     \caption{github记录Q9}
%     \label{fig:github记录Q9}
% \end{figure}

\section{第 10 题 让编程智能体进入可验证的修复循环}

\subsection{实验内容}

如图\ref{fig:复制q10}所示，使用 rsync 将 q09 复制为 q10（用 \verb|--exclude| 排除 venv、build 等构建产物），并在 q10 中创建 test\_cli.py：当 name 只含空白字符时，main 应以 \texttt{SystemExit(2)} 结束。

\begin{figure}[H]
    \centering
    \includegraphics[width=0.8\textwidth]{截屏2026-09-09 11.00.50.png}
    \caption{复制 q09 为 q10 并创建 test\_cli.py}
    \label{fig:复制q10}
\end{figure}

如图\ref{fig:测试先失败}所示，先运行 \verb|PYTHONPATH=src python -m pytest -v| 确认测试失败：报错 \texttt{TypeError: main() takes 0 positional arguments but 1 was given}。原因是当前 main 函数既不接收参数、也不会校验空白姓名，与测试期望不符，符合"先运行测试确认失败"的要求。

\begin{figure}[H]
    \centering
    \includegraphics[width=0.8\textwidth]{截屏2026-09-09 11.06.43.png}
    \caption{先运行测试确认失败}
    \label{fig:测试先失败}
\end{figure}

接下来向编程智能体说明目标（name 为纯空白时 main 应以 \texttt{SystemExit(2)} 结束）、约束（只修改实现、不修改测试）和测试命令，要求它修改实现并运行测试。如图\ref{fig:智能体修复}所示，智能体把 main 改为接收可选 argv 参数，并在 parse 后校验 \texttt{a.name.strip()}，空白时调用 p.error 使程序以退出码 2 终止；修改后再次运行测试，输出 \texttt{1 passed}。

\begin{figure}[H]
    \centering
    \includegraphics[width=0.8\textwidth]{截屏2026-09-09 11.07.39.png}
    \caption{智能体修改 cli.py 后测试通过}
    \label{fig:智能体修复}
\end{figure}

如图\ref{fig:diff人工检查}所示，人工检查 diff：\verb|diff ../q09/src/greetlab/cli.py src/greetlab/cli.py| 显示只有 main 函数签名与空白校验两处预期改动；pyproject.toml 和 \_\_init\_\_.py 与 q09 完全一致，说明智能体没有引入无关修改。

\begin{figure}[H]
    \centering
    \includegraphics[width=0.8\textwidth]{截屏2026-09-09 11.08.29.png}
    \caption{人工检查 diff 确认无无关修改}
    \label{fig:diff人工检查}
\end{figure}

如图\ref{fig:再次测试与ai_log}所示，再次运行测试全部通过；并用不超过 5 行的 ai\_log.md 记录核心提示、智能体改动和人工验证三部分内容。

\begin{figure}[H]
    \centering
    \includegraphics[width=0.8\textwidth]{截屏2026-09-09 11.09.41.png}
    \caption{再次运行测试并记录 ai\_log.md}
    \label{fig:再次测试与ai_log}
\end{figure}

\subsection{实验结果}

第十题完成了"测试先行失败、智能体修复、人工检查 diff、再次验证"的可验证修复循环：测试从 \texttt{1 failed} 变为 \texttt{1 passed}，diff 仅含预期改动，ai\_log.md 用 4 行记录完整闭环。

% TODO: 完成 w3 目录的 git 提交与推送后，将 commit/push 记录截图放入 images/ 并补充如下内容
% 第十题的 github 的 commit 记录和 push 记录如图\ref{fig:github记录Q10}所示。
% \begin{figure}[H]
%     \centering
%     \includegraphics[width=0.8\textwidth]{github记录Q10.png}
%     \caption{github记录Q10}
%     \label{fig:github记录Q10}
% \end{figure}

\section{第 11 题 把"无法处理"的协作材料改成可执行信息}

\subsection{实验内容}

围绕"空姓名仍输出问候语"这一缺陷，改写三段质量较差的协作材料。如图\ref{fig:communication初版}所示，在 q11 的 communication.md 中重写三部分内容：Issue 包含环境、复现命令、期望结果和实际结果，未知信息写"待确认"；提交信息标题使用祈使语气，正文说明问题和解决方案；评审意见指出具体行为、风险和建议动作，并标注 Blocking、Suggestion 或 Nit。

\begin{figure}[H]
    \centering
    \includegraphics[width=0.8\textwidth]{截屏2026-09-09 11.11.21.png}
    \caption{重写 communication.md 并核验字数}
    \label{fig:communication初版}
\end{figure}

经 \verb|wc -m| 核验，初版全文为 403 字，超出 400 字限制。如图\ref{fig:communication精简}所示，精简冗余表述后全文为 384 字，满足"全文控制在400字以内"的要求，且保持专业、简洁。

\begin{figure}[H]
    \centering
    \includegraphics[width=0.8\textwidth]{截屏2026-09-09 11.11.58.png}
    \caption{精简后全文 384 字}
    \label{fig:communication精简}
\end{figure}

\subsection{实验结果}

第十一题最终 communication.md 全文 384 字（$\leq$400）：Issue 部分补全了环境、复现命令、期望结果和实际结果，并将"是否仅 Windows 复现"标注为待确认；提交信息以 \texttt{fix: 空白姓名改为报错退出} 为祈使式标题，正文说明了问题与解决方案；评审意见包含 \texttt{Blocking}（未校验导致误判成功）、\texttt{Suggestion}（补充合法输入提示）和 \texttt{Nit}（校验逻辑可独立成函数便于单测）。

% TODO: 完成 w3 目录的 git 提交与推送后，将 commit/push 记录截图放入 images/ 并补充如下内容
% 第十一题的 github 的 commit 记录和 push 记录如图\ref{fig:github记录Q11}所示。
% \begin{figure}[H]
%     \centering
%     \includegraphics[width=0.8\textwidth]{github记录Q11.png}
%     \caption{github记录Q11}
%     \label{fig:github记录Q11}
% \end{figure}

\section{第 12 题 修复一个可复现的线性回归训练循环}

\subsection{实验内容}

如图\ref{fig:安装torch}所示，在 \verb|course_env| 中执行 \verb|pip install torch| 安装 PyTorch（torch 2.14.0）。

\begin{figure}[H]
    \centering
    \includegraphics[width=0.8\textwidth]{截屏2026-09-09 11.13.25.png}
    \caption{安装 PyTorch}
    \label{fig:安装torch}
\end{figure}

如图\ref{fig:创建train}所示，创建 q12/train.py，在训练循环中补全 TODO：依次调用 \verb|opt.zero_grad()| 清空梯度、\verb|loss.backward()| 反向传播、\verb|opt.step()| 更新参数；训练结束后调用 \verb|model.eval()| 进入评估模式，并在 \verb|torch.no_grad()| 中计算并打印最终损失、weight 和 bias。

\begin{figure}[H]
    \centering
    \includegraphics[width=0.8\textwidth]{截屏2026-09-09 11.14.57.png}
    \caption{创建 train.py 并补全 TODO 后运行}
    \label{fig:创建train}
\end{figure}

首次运行时 torch 提示缺少 numpy；如图\ref{fig:安装numpy重跑}所示，安装 numpy 后重新运行，输出 \texttt{final loss: 2.3923129102709773e-12}、\texttt{weight: 2.99999737739563}、\texttt{bias: -0.9999998807907104}。最终损失远小于题目要求的 0.001，weight 收敛到约 3、bias 收敛到约 $-1$，且全程没有直接把 weight 和 bias 赋值为 3 和 $-1$。

\begin{figure}[H]
    \centering
    \includegraphics[width=0.8\textwidth]{截屏2026-09-09 11.16.06.png}
    \caption{安装 numpy 后重新运行 train.py}
    \label{fig:安装numpy重跑}
\end{figure}

如图\ref{fig:目录核验}所示，最后依次查看各实验目录的文件清单，确认 q09--q12 的实验产物完整。

\begin{figure}[H]
    \centering
    \includegraphics[width=0.8\textwidth]{截屏2026-09-09 11.17.59.png}
    \caption{核验实验目录文件清单}
    \label{fig:目录核验}
\end{figure}

\subsection{实验结果}

第十二题通过补全训练循环（zero\_grad、backward、step）和评估阶段的 no\_grad 计算，完成了 PyTorch CPU 线性回归训练：最终损失 $2.39\times10^{-12}<0.001$，weight 收敛至 $2.999997\approx 3$，bias 收敛至 $-0.9999999\approx -1$，训练成功。

% TODO: 完成 w3 目录的 git 提交与推送后，将 commit/push 记录截图放入 images/ 并补充如下内容
% 第十二题的 github 的 commit 记录和 push 记录如图\ref{fig:github记录Q12}所示。
% \begin{figure}[H]
%     \centering
%     \includegraphics[width=0.8\textwidth]{github记录Q12.png}
%     \caption{github记录Q12}
%     \label{fig:github记录Q12}
% \end{figure}

\section{Github仓库地址}

本次课程的报告、练习、代码等，均已上传至个人github仓库。\\Github仓库地址：\url{https://github.com/lcx821/lcx_course_work}

